\chapter*{Foundations: Layers, Addresses, and Frames}
\addcontentsline{toc}{chapter}{Foundations: Layers, Addresses, and Frames}
\markboth{Foundations}{}
\label{ch:foundations}

Every later chapter assumes three things: that a message is wrapped in
layers, that each layer carries its own kind of address, and that a
device only reads as far up the stack as its job requires. This chapter
is the reference for all three. Nothing here is configured in eNSP ---
it is the vocabulary the labs are written in.

\section*{Two models, one stack}
OSI has seven layers and is the language people use to \emph{describe}
faults. TCP/IP has four and is what the software actually implements.
The labs use OSI numbering (``a Layer-2 problem'', ``a Layer-3 device'')
and TCP/IP behaviour.

\begin{utpfig}[OSI and TCP/IP]{The two reference models side by side. Layer numbers
come from OSI; the protocols in the right-hand column are what the eNSP labs
actually configure.}{fig:osi-tcpip}
\begin{tikzpicture}[node distance=1.2mm]
% ---- OSI column
\node[layerapp] (o7) {Application};
\node[layerapp,below=of o7] (o6) {Presentation};
\node[layerapp,below=of o6] (o5) {Session};
\node[layertra,below=of o5] (o4) {Transport};
\node[layernet,below=of o4] (o3) {Network};
\node[layerlnk,below=of o3] (o2) {Data Link};
\node[layerphy,below=of o2] (o1) {Physical};

\foreach \n/\num in {o7/7,o6/6,o5/5,o4/4,o3/3,o2/2,o1/1}
  \node[layernum,left=1.5mm of \n] {\num};

\node[font=\bfseries\small\color{UTPNavy},above=2.5mm of o7] {OSI};

% ---- TCP/IP column
\node[layerapp,right=2.6cm of o6,minimum height=26.4mm] (t4) {Application};
\node[layertra,below=of t4,minimum height=8mm] (t3) {Transport};
\node[layernet,below=of t3,minimum height=8mm] (t2) {Internet};
\node[layerlnk,below=of t2,minimum height=17.2mm] (t1) {Network Access};
\node[font=\bfseries\small\color{UTPNavy},above=2.5mm of t4] {TCP/IP};

% ---- protocols / PDU column
\node[anchor=west,font=\scriptsize\color{UTPInk},align=left,right=3mm of t4]
  {HTTP, \textbf{FTP}, \textbf{DHCP}, DNS\\[1pt]
   \textcolor{UTPTeal}{\bfseries PDU: data}};
\node[anchor=west,font=\scriptsize\color{UTPInk},align=left,right=3mm of t3]
  {TCP, UDP\\[1pt]\textcolor{UTPTeal}{\bfseries PDU: segment}};
\node[anchor=west,font=\scriptsize\color{UTPInk},align=left,right=3mm of t2]
  {IPv4, ICMP, \textbf{OSPF}\\[1pt]\textcolor{UTPTeal}{\bfseries PDU: packet}};
\node[anchor=west,font=\scriptsize\color{UTPInk},align=left,right=3mm of t1]
  {Ethernet, \textbf{802.1Q}, \textbf{PPP}, \textbf{HDLC}\\[1pt]
   \textcolor{UTPTeal}{\bfseries PDU: frame} $\rightarrow$ bits};

% ---- brackets showing the collapse
\draw[decorate,decoration={brace,amplitude=3.5pt},draw=UTPMuted]
  ($(o7.north east)+(1.5mm,0)$) -- ($(o5.south east)+(1.5mm,0)$);
\draw[decorate,decoration={brace,amplitude=3.5pt},draw=UTPMuted]
  ($(o2.north east)+(1.5mm,0)$) -- ($(o1.south east)+(1.5mm,0)$);
\end{tikzpicture}
\end{utpfig}

\begin{keyidea}
Bold entries are configured somewhere in this book. The course spends
almost all of its time at Layers 2 and 3, which is where a switch and a
router respectively make their forwarding decisions.
\end{keyidea}

\section*{Encapsulation: each layer adds its own header}
A layer treats everything handed down to it as opaque payload and puts
its own header in front. The receiver reverses the process. This is why
a single Wireshark row can be expanded into four or five nested
protocols.

\begin{utpfig}[Encapsulation]{Sending data down the stack. Each layer prepends a
header; only Ethernet also appends a trailer. The receiving host strips them in
the opposite order.}{fig:encapsulation}
\begin{tikzpicture}[x=1cm,y=1cm]
\def\u{0.30}
% row 4: application data
\protorow{0}{\u}{14/Application data/hdrpay}
\node[anchor=west,font=\scriptsize\bfseries\color{UTPPurple}] at (5.0,0) {Data};

% row 3: TCP
\protorow{-1.05}{\u}{4/TCP hdr/hdrg, 14/Application data/hdrpay}
\node[anchor=west,font=\scriptsize\bfseries\color{UTPTeal}] at (5.7,-1.05) {Segment};

% row 2: IP
\protorow{-2.10}{\u}{4/IP hdr/hdrb, 4/TCP hdr/hdrg, 14/Application data/hdrpay}
\node[anchor=west,font=\scriptsize\bfseries\color{UTPBlue}] at (6.9,-2.10) {Packet};

% row 1: Ethernet
\protorow{-3.15}{\u}{4/Eth hdr/hdro, 4/IP hdr/hdrb, 4/TCP hdr/hdrg,
                     14/Application data/hdrpay, 3/FCS/hdro}
\node[anchor=west,font=\scriptsize\bfseries\color{UTPOrange}] at (8.8,-3.15) {Frame};

% down arrow
\draw[-{Stealth[length=2.5mm]},thick,draw=UTPMuted]
  (-0.75,0.15) -- (-0.75,-3.5);
\node[rotate=90,anchor=south,font=\scriptsize\color{UTPMuted}] at (-0.95,-1.7)
  {sending host};
\end{tikzpicture}
\end{utpfig}

\section*{Three addresses, three jobs}
Beginners conflate these constantly. They answer different questions and
they change at different rates as a packet crosses a network.

\begin{center}
\begin{tabularx}{\linewidth}{>{\bfseries}p{0.15\linewidth}p{0.15\linewidth}p{0.20\linewidth}X}
\toprule
Address & Layer & Answers & Changes in transit?\\\midrule
MAC & 2 & Which device on \emph{this} link? & Yes --- rewritten at every router hop.\\
IP & 3 & Which host on \emph{any} network? & No --- constant end to end.\\
Port & 4 & Which program on that host? & No.\\\bottomrule
\end{tabularx}
\end{center}

\begin{utpfig}[Addresses across a hop]{The IP addresses stay put; the MAC addresses
are replaced on every link. A router is the device that performs the
rewrite.}{fig:addr-rewrite}
\begin{tikzpicture}[node distance=1.5cm]
\node[host] (a) {PC1};
\node[router,right=2.6cm of a] (r) {R1};
\node[host,right=2.6cm of r] (b) {PC2};
\draw[link] (a) -- (r);
\draw[link] (r) -- (b);

\node[netlabel,above=0.5mm of a] {10.0.1.10};
\node[netlabel,above=0.5mm of b] {10.0.2.20};
\node[netlabel,above=0.5mm of r] {R1 routes between them};

\node[frame,below=0.9cm of a,text width=3.0cm,align=center,draw=UTPOrange]
  {\scriptsize src MAC \textbf{aa}\\ \scriptsize dst MAC \textbf{r1}\\[1pt]
   \scriptsize src IP 10.0.1.10\\ \scriptsize dst IP 10.0.2.20};
\node[frame,below=0.9cm of b,text width=3.0cm,align=center,draw=UTPOrange]
  {\scriptsize src MAC \textbf{r1}\\ \scriptsize dst MAC \textbf{bb}\\[1pt]
   \scriptsize src IP 10.0.1.10\\ \scriptsize dst IP 10.0.2.20};

\draw[->,draw=UTPMuted,dotted] ($(a)+(0,-0.42)$) -- ($(a)+(0,-0.72)$);
\draw[->,draw=UTPMuted,dotted] ($(b)+(0,-0.42)$) -- ($(b)+(0,-0.72)$);
\node[font=\scriptsize\bfseries\color{UTPRed},below=2.35cm of r,text width=4.4cm,align=center]
  {MAC changed, IP untouched};
\end{tikzpicture}
\end{utpfig}

\section*{The Ethernet II frame}
This is the frame Wireshark shows in Lab~1 and the frame a VLAN tag is
inserted into in Lab~8.

\begin{utpfig}[Ethernet II frame]{Ethernet II framing. Widths are proportional to the
real field sizes. The preamble and SFD are generated by the hardware and are not
shown by Wireshark.}{fig:eth-frame}
\begin{tikzpicture}[x=1cm,y=1cm]
\def\u{0.40}
\protorow{0}{\u}{7/Preamble/hdrm, 1/S/hdrm, 6/Destination MAC/hdrb,
                 6/Source MAC/hdrb, 2/Type/hdro, 11/Payload/hdrpay, 4/FCS/hdrg}
\prototicks{-0.63}{\u}{7/7, 1/1, 6/6, 6/6, 2/2, 11/46--1500, 4/4}
\node[hdrtick,anchor=east] at (-0.12,-0.63) {bytes};
\draw[decorate,decoration={brace,amplitude=3pt,mirror},draw=UTPMuted]
  (0,-0.95) -- (3.2,-0.95);
\node[hdrtick] at (1.6,-1.30) {added by the NIC};
\draw[decorate,decoration={brace,amplitude=3pt,mirror},draw=UTPTeal]
  (3.25,-0.95) -- (14.75,-0.95);
\node[hdrtick,text=UTPTeal] at (9.0,-1.30) {what a capture shows};
\end{tikzpicture}
\end{utpfig}

\begin{keyidea}
\cmd{Type} tells the receiver which Layer-3 protocol is inside:
\cmd{0x0800} for IPv4, \cmd{0x0806} for ARP, \cmd{0x8100} for a
VLAN-tagged frame. Wireshark reads this field to decide which dissector
to run, which is why an unrecognised type shows as raw data.
\end{keyidea}

\section*{The IPv4 header}
Four fields matter for these labs: the two addresses, the TTL that stops
routing loops, and the protocol number that identifies the payload.

\begin{utpfig}[IPv4 header]{The IPv4 header, drawn as the 32-bit rows Wireshark
displays. Shaded fields are the ones the labs read or change.}{fig:ipv4-header}
\begin{tikzpicture}[x=1cm,y=1cm]
\def\u{0.36}
\protorow{0}{\u}{4/Ver/hdrm, 4/IHL/hdrm, 8/DSCP + ECN/hdrm, 16/Total length/hdrm}
\protorow{-0.85}{\u}{16/Identification/hdrm, 3/Flg/hdrm, 13/Fragment offset/hdrm}
\protorow{-1.70}{\u}{8/TTL/hdro, 8/Protocol/hdro, 16/Header checksum/hdrm}
\protorow{-2.55}{\u}{32/Source IP address/hdrb}
\protorow{-3.40}{\u}{32/Destination IP address/hdrb}

\node[hdrtick,anchor=west] at (0,0.62) {0};
\node[hdrtick] at (5.76,0.62) {16};
\node[hdrtick,anchor=east] at (11.52,0.62) {31};
\draw[draw=UTPMuted!50] (0,0.45)--(11.52,0.45);

\node[anchor=west,font=\scriptsize\color{UTPMuted},align=left] at (11.9,-1.70)
  {\textcolor{UTPOrange}{\bfseries TTL} drops by 1 per hop.\\
   \textcolor{UTPOrange}{\bfseries Protocol}: 1 ICMP,\\ 6 TCP, 17 UDP, 89 OSPF.};
\node[anchor=west,font=\scriptsize\color{UTPMuted},align=left] at (11.9,-2.98)
  {\textcolor{UTPBlue}{\bfseries The pair a router}\\
   \textcolor{UTPBlue}{\bfseries matches against its}\\
   \textcolor{UTPBlue}{\bfseries routing table.}};
\end{tikzpicture}
\end{utpfig}

\section*{What each device reads}
A device only decodes as far up as it needs to. This single fact
explains why a switch cannot fix a routing problem and why a VLAN
cannot be bypassed by changing a PC's IP address.

\begin{utpfig}[Device scope]{How far up the stack each device looks. The dividing
line between a switch and a router is the whole of Chapters~\ref{ch:static-routing}
and~\ref{ch:ospf}.}{fig:device-scope}
\begin{tikzpicture}[x=1cm,y=1cm]
\foreach \x/\name/\desc in {%
  0/Hub/{repeats bits},%
  3.5/Switch/{learns MAC},%
  7.0/Router/{reads IP},%
  10.5/Firewall/{reads ports}}
  {\node[font=\bfseries\small\color{UTPNavy}] at (\x+0.9,0.55) {\name};
   \node[font=\scriptsize\color{UTPMuted}] at (\x+0.9,0.15) {\desc};}

% layer bands
\foreach \y/\lab/\st in {-0.4/Physical/hdrm, -1.1/Data Link/hdro,
                         -1.8/Network/hdrb, -2.5/Transport/hdrg}
  {\node[\st,minimum width=1.6cm,anchor=east,font=\scriptsize\color{UTPNavy}]
     at (-0.25,\y) {\lab};}

\foreach \x/\top in {0/1, 3.5/2, 7.0/3, 10.5/4} {
  \foreach \i in {1,...,4} {
    \pgfmathsetmacro{\yy}{-0.4-(\i-1)*0.7}
    \ifnum\i>\top
      \node[draw=UTPMuted!30,fill=white,minimum width=1.8cm,minimum height=5.5mm,
            font=\tiny\color{UTPMuted!60}] at (\x+0.9,\yy) {--};
    \else
      \node[draw=UTPGreen,fill=UTPPaleTeal,minimum width=1.8cm,minimum height=5.5mm,
            font=\tiny\bfseries\color{UTPGreen}] at (\x+0.9,\yy) {reads};
    \fi
  }
}
\end{tikzpicture}
\end{utpfig}

\begin{keyidea}
A switch forwards on MAC addresses and never looks at an IP address.
Two PCs in different VLANs stay unreachable however you renumber them,
because the switch stops reading before it reaches the IP header.
\end{keyidea}

\section*{Collision and broadcast domains}
The terms recur in Chapters~\ref{ch:lacp} and~\ref{ch:vlan}. A hub gives
one collision domain and one broadcast domain. A switch gives one
collision domain per port but still one broadcast domain. Only a router
or a VLAN splits the broadcast domain.

\begin{utpfig}[Domains]{What each device splits. The middle case is the one people
get wrong: a switch gives every port its own collision domain but leaves the
broadcast domain intact.}{fig:domains}
\begin{tikzpicture}[x=1cm,y=1cm,every node/.style={font=\scriptsize}]
\foreach \x/\dev/\coll/\bcast in {0/Hub/1/1, 5.0/Switch/4/1, 10.0/Router/4/4} {
  \node[font=\bfseries\small\color{UTPNavy}] at (\x+1.6,1.5) {\dev};
  \node[switch,minimum width=2.6cm] at (\x+1.6,0.75) {};
  \foreach \i in {0,1,2,3} {
    \node[host,minimum width=0.62cm,minimum height=0.42cm,font=\tiny]
      at (\x+0.45+\i*0.77,-0.35) {};
    \draw[link] (\x+0.45+\i*0.77,-0.15) -- (\x+0.45+\i*0.77,0.45);
  }
  \node[align=center,color=UTPMuted] at (\x+1.6,-1.0)
    {\textbf{\coll} collision \\ \textbf{\bcast} broadcast};
}
% collision domain shading
\begin{scope}[on background layer]
\draw[bcastzone] (-0.05,-0.62) rectangle (3.25,0.42);
\foreach \i in {0,1,2,3}
  \draw[bcastzone] (5.15+\i*0.77,-0.62) rectangle (5.75+\i*0.77,0.42);
\foreach \i in {0,1,2,3}
  \draw[bcastalt] (10.15+\i*0.77,-0.62) rectangle (10.75+\i*0.77,0.42);
\end{scope}
\end{tikzpicture}
\end{utpfig}

\begin{quickcheck}
\begin{enumerate}
\item A packet crosses two routers. How many times is its destination
      MAC address rewritten, and how many times its destination IP?
\item Which header field would you check first to diagnose a routing loop?
\item A frame arrives with Type \cmd{0x8100}. What does the switch expect
      to find next, and which chapter configures it?
\item Why can a switch not resolve a problem caused by a missing static route?
\end{enumerate}
\end{quickcheck}
